% ======================================================================
\chapter{RL congestion control for the wide-area Internet}
\label{ch:internet}
% ======================================================================

The Internet is where RL-based congestion control started and where most of it is still done. This chapter covers the \nChInternetN{} studies whose setting is the wide-area Internet (\nSetWan) or a generic wired bottleneck (\nSetGen). Of these, \nChInternetDesign{} are design studies and \nChInternetAnalysis{} analyze existing controllers or provide frameworks for them. We first follow the main design lineage, from clean-slate learned controllers (\cref{sec:clean-slate}) through hybrids that learn on top of classic algorithms (\cref{sec:hybrid}) to data-driven and offline designs (\cref{sec:offline}). We then turn to multi-objective control (\cref{sec:multiobjective}), fairness and multi-agent learning (\cref{sec:fairness}), multipath transport (\cref{sec:multipath}), real-time media (\cref{sec:rtc}), other transports (\cref{sec:other-transports}), and network-side control (\cref{sec:aqm}). \Cref{tab:landmarks} summarizes the designs that later work most often builds on or compares against.

\begin{table}[t]
\centering
\caption{Landmark Internet designs: where the agent acts, what it observes, what it learns from, and the most realistic environment in which it was evaluated. Details are from the full texts except where marked $^\dagger$ (abstract and reviewers' knowledge). NR: not reported in the material we could access. Level: highest evaluation environment.}
\label{tab:landmarks}
\scriptsize
\setlength{\tabcolsep}{3pt}
\begin{tabularx}{\linewidth}{@{}L{2.0cm}L{2.75cm}YL{2.55cm}L{1.75cm}l@{}}
\toprule
Study & Action & Observation & Learning signal & Algorithm & Level \\
\midrule
Indigo~\cite{yan2018pantheon} & cwnd step every 10\,ms & EWMAs of queueing delay, sending and receiving rate; current cwnd; previous action & imitation of an oracle that knows the emulated link's ideal cwnd & IL, LSTM & Real \\
Aurora~\cite{jay2019aurora} & rate multiplier per MI & history of latency gradient, latency ratio, sending ratio & $10\,\text{thr} - 1000\,\text{lat} - 2000\,\text{loss}$ & PPO & Emu \\
MVFST-RL~\cite{sivakumar2019mvfst} & cwnd $\in\{\cdot, /2, -10, +10, \times 2\}$ & ACK statistics over 100\,ms plus action history & throughput--delay trade-off & IMPALA (async.) & Emu \\
Orca~\cite{abbasloo2020orca} & cwnd $= 2^{2a}\cdot$ CUBIC's cwnd, once per monitoring period & throughput, delay, loss statistics & throughput minus loss penalty, scaled by a delay factor & TD3 & Real \\
Eagle$^\dagger$~\cite{emara2020eagle} & rate & network statistics with LSTM memory & reward plus guidance from BBR as an expert & DRL + IL & Emu \\
Owl~\cite{sacco2021owl} & cwnd (Linux kernel module) & end-to-end statistics plus in-network signals & a utility function & DQN & Emu \\
MOCC~\cite{ma2022mocc} & rate & network statistics plus objective weights & preference-weighted throughput, delay, loss & PPO + preference sub-net & Real \\
Genet~\cite{xia2022genet} & (training method) & Aurora's & Aurora's & PPO + curriculum & Real \\
Sage$^\dagger$~\cite{yen2023sage} & cwnd & a general representation of many path statistics & logged behavior of existing heuristics & offline RL & Real \\
Spine$^\dagger$~\cite{tian2024spine} & sub-policy for a light CC executor & NR & NR & hierarchical DRL & Real \\
Astraea~\cite{liao2024astraea} & rate & local, global (multi-flow) and historical statistics & throughput, delay, loss, fairness & multi-agent actor-critic & Real \\
Canopy~\cite{yang2024canopy} & as Orca & as Orca & Orca's reward plus verifier feedback on properties & TD3 + certificates & Real \\
Mutant~\cite{pappone2025mutant} & select a kernel CC scheme & raw transport statistics & power (throughput over delay) & LinUCB bandit & Real \\
\bottomrule
\end{tabularx}
\end{table}

% ----------------------------------------------------------------------
\section{Clean-slate learned controllers}
\label{sec:clean-slate}

\paragraph{Q-learning beginnings.} The earliest Internet studies in the corpus use Q-learning with simple function approximation. QTCP~\cite{li2019qtcp,li2017dynamic} approximates the Q-function with (generalized) Kanerva coding and reports 59.5\% higher throughput than rule-based TCP. RL-TCP~\cite{kong2018improving} learns cwnd adjustments for shallow buffers in ns-2. These designs showed that a learned cwnd policy could beat NewReno in simulation. However, their state spaces were small and hand-quantized, and none reports an evaluation beyond simulation. Later work showed that QTCP flows are selfish when they compete with each other. Wrapping the learned decision in an AIMD rule restores fairness without losing throughput~\cite{yamazaki2021fairness}.

\paragraph{Aurora and the deep-RL formulation.} Aurora~\cite{jay2019aurora} established the formulation most later work adopts (\cref{sec:formulation}): PPO, monitor intervals, a history of scale-free features, and a reward that trades throughput against latency and loss. Trained on simulated links of 1.2--6\,Mbit/s with 50--500\,ms latency and 0--5\% random loss, it kept near-full utilization in emulation from 1 to 128\,Mbit/s, more than 20$\times$ its training range, and over wide ranges of buffer size and loss, comparing well with CUBIC and PCC Vivace. It performed poorly only at 1\,ms latency, where about 1\,ms of packet-processing noise that was absent from training dominated the latency signal. Its authors explicitly flagged fairness, safety and generalization as open problems. Seven years later, those three problems are still the main research agenda (\cref{ch:training,ch:deployment}). Aurora's open code made it the default testbed for methods: curriculum learning~\cite{xia2022genet}, teacher--student training~\cite{zheng2022teacher}, scenario-based reward shaping~\cite{yerushalmi2023enhancing}, interpretation~\cite{meng2020interpreting}, symbolic distillation~\cite{sharan2022symbolic}, formal verification~\cite{kazak2019verily,eliyahu2021whirl,amir2024verifying}, meta-RL~\cite{mwaura2025adaptive} and fairness shaping~\cite{mbrice2026fair} have all been demonstrated on Aurora.

\paragraph{Imitation learning and training at scale.} In parallel, Indigo~\cite{yan2018pantheon}, released with the Pantheon testbed, trained an LSTM by imitation learning. Its oracle knew the ideal congestion window of each emulated link, and Pantheon's emulators were calibrated against real paths so that this training would transfer. MVFST-RL~\cite{sivakumar2019mvfst} addressed a systems problem that synchronous simulators hide: in a real stack, the network does not pause while the agent computes. It runs the policy asynchronously inside Facebook's QUIC implementation and trains with IMPALA's off-policy correction. Other studies learn online or directly on real networks. ARC~\cite{zhang2020reinforcement} trains asynchronously in real networks, Rax~\cite{bachl2019rax} learns online within minutes, and MARLIN~\cite{galliera2023marlin} trains SAC on a real network and then specializes it for tactical satellite and UHF links.

\paragraph{The broad middle of the literature.} Many Internet design studies follow the Aurora template with a different algorithm, feature set or network architecture. Examples are DQN-family controllers such as TCP-Drinc~\cite{xiao2019drinc}, RLCC~\cite{xia2021rlcc}, TCP-DQN~\cite{wang2021intelligent}, TCP-QNCC~\cite{huang2023tcp}, KerDqn~\cite{xu2024kerdqn} and others~\cite{aglamazlar2025deep,aoki2026performance}; PPO controllers~\cite{li2020tcp,huang2022tcp,shi2023intelligent,yuan2023research}; and SAC or other actor-critic controllers~\cite{guo2022deep,xing2025reinforcement,blokhin2024deep}. Typical claims are double-digit throughput or delay gains over NewReno or CUBIC on a few scenarios. Such gains are difficult to compare, because the scenarios, baselines and metrics differ from study to study (\cref{ch:training}). A smaller group studies the formulation itself, and its findings are more transferable than any single controller. The length of the monitor interval and the RTT estimator change QoS noticeably~\cite{naqvi2023fine}. Discrete actions learn more easily than continuous ones, and fourteen candidate rewards produce visibly different operating points~\cite{kiarie2025design}. Rainbow's DQN extensions do not behave in congestion control as they do in Atari~\cite{martins2021rainbow}. Normalizing the reward by an estimated throughput ceiling makes it scale-invariant, so that one policy can serve links of different speeds~\cite{chae2025ackurate,chae2026pacemaker}. Explicit probing mechanisms or intrinsic probing rewards help an agent discover spare bandwidth~\cite{li2021tcp,han2026learn}.

% ----------------------------------------------------------------------
\section{Hybrid designs: learning on top of classic control}
\label{sec:hybrid}

Hybrid designs are the clearest trend in the corpus. They grow from \nPctHybridPre\% of design studies before 2023 to \nPctHybridPost\% since, and \nChInternetHybrid{} of the \nChInternetDesign{} Internet design studies are hybrid. They assume that a neural policy acting at packet timescales is too slow, too expensive and too unpredictable, while classic algorithms react quickly and safely but cannot adapt their operating point. The designs differ in how they divide the work.

\paragraph{The agent scales a classic window.} Orca~\cite{abbasloo2020orca} is the reference design. CUBIC runs per ACK in the kernel, and once per monitoring period a TD3 agent multiplies CUBIC's window by $2^{2a}$, $a\in[-1,1]$. Orca was evaluated on emulated links and on global testbeds and found to be more consistent than clean-slate learned controllers at a fraction of their overhead. Canopy~\cite{yang2024canopy} keeps Orca's architecture but adds \emph{quantitative certificates} from a verifier to the reward. The verifier measures how close the policy is to satisfying operator-specified safety and robustness properties, such as robustness to measurement noise. This raises worst-case property satisfaction by up to 1.4$\times$ while preserving or improving average performance. Other designs build the agent into BBR: ABBR infers BBR's pacing rate with DRL while staying compatible with vanilla BBR~\cite{pan2023making}, and Marten guides a DRL agent with BBR and shields it at run time~\cite{zhou2025safety}. Eagle uses BBR differently, as a teacher during training only~\cite{emara2020eagle}. PBQ splits the knobs: PPO sets the window while BBR sets the pacing rate~\cite{zhang2023pbq}. HAN falls back to a classic controller when the policy's uncertainty is high~\cite{jia2026han}.

\paragraph{The agent tunes parameters.} A lighter variant leaves the control law intact and tunes its parameters: CUBIC's $\beta$ and $C$~\cite{sharma2026adaptive}, CUBIC's window thresholds~\cite{abedini2025cubic}, LEDBAT++'s delay target~\cite{tang2023combining}, or several explicit-feedback parameters~\cite{kheirkhah2023acct}. Spine~\cite{tian2024spine} generalizes the idea hierarchically. A coarse-grained DRL agent emits \emph{sub-policies} for a lightweight executor that reacts at fine granularity, which cuts control overhead by at least 80\% relative to DRL controllers and makes it comparable to CUBIC.

\paragraph{Classic and learned candidates compete.} In Libra~\cite{du2025libra,du2021libra}, a classic algorithm and an RL policy each propose a rate. Libra briefly tries both candidates in an evaluation stage and commits to whichever achieves the higher utility. It reports 1.2$\times$ Orca's throughput on dynamic networks with good fairness and convergence. Anole~\cite{han2025anole} extends this with candidate selection and deprecation.

\paragraph{The agent selects the algorithm.} The coarsest hybrid learns \emph{which} existing algorithm to run. TCP-RL~\cite{nie2019tcprl}, deployed at a large search engine, combines group-based RL for the initial window with deep RL that selects among 14 algorithms. Floo~\cite{zhang2023floo} selects algorithms for a production mobile web service and switches between them smoothly. Mutant~\cite{pappone2025mutant} uses a LinUCB contextual bandit to switch among a pool of kernel protocols in real time, with negligible harm to competing flows. Similar selectors use offline-trained DRL~\cite{li2026alchemy}, contextual Thompson sampling~\cite{peng2025htpccs} or meta-evolutionary RL~\cite{kanungo2026adaptive}.

\finding{Hybrid designs trade some of the theoretical flexibility of a clean-slate policy for the reaction speed, low overhead and predictable failure behavior of classic algorithms. The Internet controllers in the corpus that report production use are either hybrids with a classic algorithm or safeguard in the loop (TCP-RL, Loki, Floo, Marten, and OnRL in \cref{ch:wireless}) or policies imitated from logs or an oracle (Concerto, DuGu). None is a clean-slate policy trained only by exploring a simulator (\cref{ch:deployment}).}

% ----------------------------------------------------------------------
\section{Data-driven and offline designs}
\label{sec:offline}

Online RL needs an environment to explore, and for congestion control that environment is almost always a simulator or emulator, which is the source of the generalization problem. A newer line learns from \emph{logged} data. Sage~\cite{yen2023sage} trains offline on trajectories produced by a pool of existing heuristic schemes. Its authors present it as the first purely data-driven Internet design, and in Internet and controlled tests it learned a policy that performs better than the heuristics it observed. POCC~\cite{jia2026pocc} adds pessimism (critic ensembles) to guard against extrapolation from multi-algorithm data. A decision-transformer controller~\cite{ji2026dtcc} treats control as sequence modelling. Earlier, DuGu~\cite{huang2022learned} trained a network to imitate an optimal solver running in an omniscient emulator and validated it with real-world A/B tests. Offline learning also dominates real-time media, where operators have logs from millions of calls (\cref{sec:rtc}).

% ----------------------------------------------------------------------
\section{Multiple objectives and application preferences}
\label{sec:multiobjective}

A single reward hard-codes one operating point, yet applications need different ones: throughput for bulk transfers, low delay for interactive applications, and yielding for scavenger traffic. MOCC~\cite{ma2022mocc} conditions a PPO policy on a preference vector. One model then competes with the best specialized algorithm for each objective, and transfer learning lets it adapt to new objectives 14.2$\times$ faster than prior work. Zhang et al.\ train a family of near-Pareto policies with DDPG and select one per weighting without retraining~\cite{zhang2022deepcc}, and other studies use constrained RL with Lagrangian weights~\cite{liu2021multi}, multi-objective RL distilled into a decision tree~\cite{liu2025adaptive}, or a client--server architecture that accepts arbitrary application objectives~\cite{xing2025reinforcement}. Deadline-aware variants~\cite{liao2021learning,liao2023stop} and a learned less-than-best-effort controller~\cite{tang2023combining} target specific application classes. Reward design is also studied directly. Zygos corrects the reward distribution with meta-gradients~\cite{liu2024zygos}, and scenario-based modelling lets engineers inject rules into the reward~\cite{yerushalmi2023enhancing}.

% ----------------------------------------------------------------------
\section{Fairness and multi-agent learning}
\label{sec:fairness}

A policy trained alone on a link learns to be greedy. When several copies share a bottleneck, nothing in a single-agent reward pushes them toward an equal share. Fairness is evaluated in \nChInternetFair{} of the \nChInternetDesign{} Internet design studies, and the corpus offers three remedies. The first is \emph{multi-agent training}. Pareto trains communication-free agents in stages and refines them online~\cite{emara2022pareto}. Astraea trains flows jointly in a ``faithful'' multi-flow environment with a fairness term in the reward, reaching near-optimal fairness with up to 8.4$\times$ faster convergence~\cite{liao2024astraea}. MAPPO-based~\cite{peng2025macc} and meta-RL~\cite{bai2025pisces} variants follow. The second is to \emph{decouple} fairness from the learned policy. Jury transforms network signals into a view that is common to all competing flows, and then post-processes the policy's sending rate according to each flow's estimated share of the bandwidth. Large flows thus behave more conservatively and small flows more aggressively, and fairness holds even in conditions unseen in training~\cite{tian2025jury}. The third is \emph{theory}. Learning algorithms with swap-regret guarantees provably converge to correlated equilibria among competing senders~\cite{huang2023end,huang2024game}. Independent evaluations are less favorable. Giacomoni and Parisis find public RL controllers unfair in almost all setups~\cite{giacomoni2024rlcc}, and Mazilu et al.\ show that a fairness term in the reward works in training conditions but does not generalize~\cite{mazilu2025learning}. Fairness against the \emph{deployed} base (CUBIC, BBR) is evaluated far less often than fairness among copies of the same learned controller.

% ----------------------------------------------------------------------
\section{Multipath transport}
\label{sec:multipath}

Multipath TCP and QUIC add a second decision, namely how to split traffic over subflows, and a coupling constraint, namely not to take more than a single-path flow at a shared bottleneck. Learned multipath controllers appear in \nTopicMP{} studies across all settings (wireless ones in \cref{ch:wireless}). On generic wired paths, DRL-CC~\cite{xu2019drlcc} uses an LSTM actor-critic over all subflows and remains TCP-friendly. Partner~\cite{huang2020objective} aligns the controller's reward with the scheduler's. DQL-MPTCP~\cite{pokhrel2021learning} combines deep Q-learning with policy gradients and an LSTM, and DDPG-MPCC~\cite{pokhrel2024ddpg} combines DDPG with online convex optimization in a kernel implementation. Bandit designs are common too: MP-OL runs a bandit per subflow that can fall back to a traditional controller~\cite{zhuang2022achieving}, and MPLibra brings Libra's classic/learned competition to multipath~\cite{yu2024classic}. Maggey trains a single-subflow model and composes it with non-blocking mode transitions~\cite{wei2024breaking}.

% ----------------------------------------------------------------------
\section{Real-time media}
\label{sec:rtc}

Interactive video (video calls, cloud gaming, live streaming) is the application area where learned congestion control has been deployed most. \nTopicRTC{} studies target it. The controller is a \emph{bandwidth estimator} whose output sets the encoder's target bitrate, replacing or tuning WebRTC's Google Congestion Control (GCC). In many designs it runs at the receiver, which estimates the bandwidth and feeds it back to the sender. Rewards are QoE terms such as video quality, frame delay and stalls.

\paragraph{From simulation to production.} QARC~\cite{huang2018qarc} paired A3C with a video-quality predictor in trace-driven simulation. Fang et al.\ at Microsoft~\cite{fang2019rtc} trained PPO for real calls and documented the gap between simulated and real reward. Concerto~\cite{zhou2019concerto} used imitation learning on logged sessions and was incorporated into a production video-telephony application. OnRL~\cite{zhang2020onrl} moved learning online. Each user runs an agent, the agents' experience is aggregated into one model, and a rule-based safeguard takes over when the policy misbehaves. In a month-long production evaluation it reduced stalling by 14.2\% at similar quality. Loki~\cite{zhang2021loki} fuses the RL model with a reverse-engineered rule-based model and was evaluated in a commercial RTC service. DuGu~\cite{huang2022learned} validated imitation-learned control of video upload with A/B tests.

\paragraph{Hybrids and benchmarks.} Several later designs are hybrids around GCC. HRCC~\cite{wang2021hybrid} lets an agent set a gain on GCC's estimate, BoB warm-starts PPO with a heuristic~\cite{bentaleb2022bob}, HybridRTS switches between rule-based and RL control~\cite{zhang2022hybridrts}, and DynA learns inside GCC with a curriculum~\cite{wu2026dyna}. Shared infrastructure has made this subfield unusually comparable. OpenNetLab provides a gym environment plus distributed real nodes~\cite{eo2022opennetlab}. ReCoCo trains on 162 traces and tests on more than 3,000 unseen ones~\cite{markudova2026advancing}, and GenRL-CC adds GAN-synthesized traces to it~\cite{islam2026genrl}. The MMSys'24 grand challenge~\cite{khairy2024acm} released logged Microsoft Teams data for \emph{offline} RL, and the best entries matched the strongest behavior policy in the logs.

\paragraph{Estimation accuracy.} Song et al.~\cite{song2025debunking} showed that a two-stage regressor, trained by imitation and then refined online, outperforms several RL bandwidth estimators. They conclude that the accuracy of the bandwidth estimate, not reward-driven learning, determines QoE. Together with the offline-RL challenge results, their finding suggests that much of the benefit in real-time media comes from learning from data rather than from RL's exploration.

% ----------------------------------------------------------------------
\section{QUIC, NDN, SDN and other transports}
\label{sec:other-transports}

QUIC's user-space implementation makes it a convenient host for learned controllers. \nTopicQUIC{} studies use it, from MVFST-RL~\cite{sivakumar2019mvfst} to SAC-based control~\cite{blokhin2024deep}, PPO with BBR pacing~\cite{zhang2023pbq}, multicast QUIC~\cite{chen2023congestion}, XQUIC~\cite{yu2025dnccq}, and video-oriented controllers that report VMAF and rebuffering gains~\cite{sidhu2026pluggable,suarez2025pulsequic}. Named Data Networking (NDN) changes the control point. Consumers pace Interests, and routers see per-content state, which learned controllers exploit at the receiver~\cite{lan2019deep,zheng2024receiver} or with edge-aided congestion degrees~\cite{yang2022ieacc}. In software-defined networks a central controller can allocate rates directly, either with tabular or deep RL~\cite{zhang2021research,lei2020congestion,boussaoud2025dual} or with multi-agent cooperative Q-learning~\cite{boussaoud2023multi}. A central agent can also periodically guide end-to-end controllers instead of replacing them~\cite{konig2026augmenting}.

% ----------------------------------------------------------------------
\section{Network-side control: AQM and ECN}
\label{sec:aqm}

\nTopicAQM{} studies learn an \emph{active queue management} policy at the router instead of the sender's control law. They range from early actor-critic and Q-learning AQMs~\cite{xu2018network,gomez2019intelligent,jiang2020performance} to DQN- and DDPG-based AQMs, some of which tune the parameters of RED or BLUE~\cite{yousif2022intelligent,matrood2024adaptive,mohammed2026ddpg}. Some select the queue discipline itself, from FIFO, WFQ, RED and CoDel~\cite{liu2026q}. Network-side learning has two practical advantages: it needs no end-host changes, and it can see the queue directly. Two studies exploit this with modern data planes. An A3C agent tunes the L4S DualQ coupled AQM in FreeBSD~\cite{satish2024active}, and DESiRED tunes a P4 AQM using in-band telemetry~\cite{almeida2024desired}. LFQ learns per-flow buffer sizes for fair queueing~\cite{bachl2020lfq}. MACC trains a transport agent and an AQM agent to cooperate~\cite{bai2022macc}. Network-side learning is far more common in data centers, where operators control the switches (\cref{ch:datacenter}).

% ----------------------------------------------------------------------
\section{Summary}
\label{sec:internet-summary}

\begin{itemize}
  \item \textbf{Design (\RQ{1}).} Internet designs act on the rate (\nChInternetActRate) or the window (\nChInternetActCwnd), usually from a short history of throughput, delay and loss statistics, with a weighted throughput--delay--loss reward. The field has moved from clean-slate policies to hybrids that delegate fast reaction to a classic algorithm, and from purely simulated exploration toward learning from logs.
  \item \textbf{Evidence (\RQ{2}).} Of the \nChInternetDesign{} Internet design studies, \nChInternetReal{} reach a real network, \nChInternetEmu{} stop at emulation, \nChInternetSim{} at simulation, and \nChInternetNR{} do not report their environment in the material we could access. Only \nChInternetGen{} explicitly test outside their training distribution, and \nChInternetCode{} release code.
  \item \textbf{Deployment (\RQ{3}).} Production use is concentrated in real-time media and in algorithm selection or parameter tuning (TCP-RL, Concerto, OnRL, Loki, Floo, DuGu). Most of these keep a classic algorithm or a safeguard in the loop; the two that do not (Concerto, DuGu) were imitated from logs or from an optimal solver rather than learned by exploration.
\end{itemize}
